% PAGES: 169
% Continues the \begin{proof} of Lemma 5.6 left open at the end of
% sec05_c2_3.tex (page 168 / folio 152), right after equation (5.58). No
% fresh "\begin{proof}" is opened here; this file closes it with
% \end{proof} where the printed square appears, after equation (5.62).

see (10.6) and (10.8) for deriving (5.57) and (5.58).

From (5.56--5.58), it follows that requiring all the leading principal minors of the
matrix $A$ to be positive is equivalent to (5.53). Then, from Sylvester's Criterion
(Theorem 5.8), we conclude that $A$ is symmetric positive definite if and only if
(5.53) is satisfied.

% Source wording: "The principal minors the matrix A are" -- missing "of"
% before "the matrix" -- reproduced as printed, per the never-correct-the-
% author rule.
(ii) The principal minors the matrix $A$ are
\begin{equation}
\det(d_1) = d_1; \quad \det(d_2) = d_2; \quad \det(d_3) = d_3;
\end{equation}
\begin{equation}
\det\begin{pmatrix} d_1 & a\\ a & d_2 \end{pmatrix} = d_1d_2 - a^2; \quad
\det\begin{pmatrix} d_1 & b\\ b & d_3 \end{pmatrix} = d_1d_3 - b^2;
\end{equation}
\begin{equation}
\det\begin{pmatrix} d_2 & c\\ c & d_3 \end{pmatrix} = d_2d_3 - c^2;
\end{equation}
\begin{equation}
\det\begin{pmatrix} d_1 & a & b\\ a & d_2 & c\\ b & c & d_3 \end{pmatrix} = d_1d_2d_3 + 2abc - d_3a^2 - d_2b^2 - d_1c^2;
\end{equation}

From (5.59--5.62), it follows that requiring all the principal minors of the matrix $A$
to be positive is equivalent to (5.54--5.55). Then, from Sylvester's Criterion (Theorem
5.8), we conclude that $A$ is symmetric positive semidefinite if and only if the
inequalities (5.54--5.55) are satisfied.
\end{proof}

Note that, by letting $d_1 = d_2 = d_3 = 1$ in (5.53), it follows that
\begin{equation}
\begin{split}
A \;&=\; \begin{pmatrix} 1 & a & b\\ a & 1 & c\\ b & c & 1 \end{pmatrix} \ \text{spd}\\
&\Longleftrightarrow \quad a^2 < 1 \ \text{and} \ \det(A) = 1+2abc-a^2-b^2-c^2 > 0\\
&\Longleftrightarrow \quad -1 < a < 1 \ \text{and} \ \det(A) = 1+2abc-a^2-b^2-c^2 > 0.
\end{split}
\end{equation}

Similarly, by letting $d_1 = d_2 = d_3 = 1$ in (5.54--5.55), it follows that
\begin{equation}
\begin{split}
A \;&=\; \begin{pmatrix} 1 & a & b\\ a & 1 & c\\ b & c & 1 \end{pmatrix} \ \text{spsd}\\
&\Longleftrightarrow \quad a^2 \leq 1;\ b^2 \leq 1;\ c^2 \leq 1;\ \det(A) = 1+2abc-a^2-b^2-c^2 \geq 0\\
&\Longleftrightarrow \quad -1 \leq a,b,c \leq 1; \ \text{and} \ \det(A) = 1+2abc-a^2-b^2-c^2 \geq 0.
\end{split}
\end{equation}

Note that the matrix $\begin{pmatrix} 1 & a & b\\ a & 1 & c\\ b & c & 1 \end{pmatrix}$ is
the correlation matrix of three random variables $X_1$, $X_2$, and $X_3$ with
correlations $\corr(X_1,X_2) = a$, $\corr(X_1,X_3) = b$, and $\corr(X_2,X_3) = c$, with
$-1 \leq a,b,c \leq 1$. The positive definiteness criteria above will be further used
in Section 7.4 to establish possible values for $a$, $b$, and $c$ such that the matrix
\[
\begin{pmatrix} 1 & a & b\\ a & 1 & c\\ b & c & 1 \end{pmatrix}
\]
is, indeed, a correlation matrix.
% Page 169 (folio 153) ends here; the sentence and paragraph appear
% complete on the source page (no environment left open). c3's first page
% (170) should be checked for whether this paragraph or section continues.
